\documentclass[11pt,a4paper]{article}
\usepackage[utf8]{inputenc}
\usepackage[T1]{fontenc}
\usepackage[english]{babel}
\usepackage{geometry}
\usepackage{amsmath,amssymb,amsfonts}
\usepackage{graphicx}
\usepackage{booktabs}
\usepackage{siunitx}
\usepackage{hyperref}
\usepackage{xcolor}
\usepackage{enumitem}
\geometry{margin=2.5cm}
\hypersetup{colorlinks=true,linkcolor=blue,citecolor=blue,urlcolor=blue}
\title{Improving a PINN for Violent Sloshing Simulation\\
\large Interface Representation, Adaptive Sampling, and Temporal Training}
\author{Ludovic Andrieux}
\date{\today}

\begin{document}
\maketitle

\begin{abstract}
[Abstract de 150--250 mots : contexte, difficulté principale, méthode proposée, protocole de validation et résultats quantitatifs à compléter.]
\par\noindent\textbf{Keywords :} PINN ; ballottement ; écoulement diphasique ; fraction volumique ; interface libre ; échantillonnage adaptatif ; optimisation.
\end{abstract}

\section{Introduction}
\subsection{Background and motivation}
Introduce tank sloshing and free-surface flows, particularly violent impacts and wave breaking.
\subsection{Limitations of numerical methods and PINNs}
Briefly describe conventional CFD methods (including interFoam) and PINN challenges: steep gradients, diffuse interfaces, imbalanced losses, and global-in-time training.
\subsection{Objectives and contributions}
List only contributions demonstrated in the paper. For example:
\begin{itemize}[nosep]
  \item encodage de Fourier optionnel des coordonnées ;
  \item échantillonnage guidé par l'interface et les résidus physiques ;
  \item entraînement par fenêtres temporelles avec transfert des poids ;
  \item pondération adaptative des pertes et optimisation Adam puis L-BFGS ;
  \item évaluation quantitative de la précision, du coût et de la localisation de l'interface.
\end{itemize}

\section{Physical Problem}
\subsection{Tank geometry and excitation}
Define the geometry, imposed motion, initial conditions, and physical parameters.
\subsection{Governing equations}
Present the incompressible flow and momentum equations with volume-fraction-dependent properties. State assumptions and the surface-tension model actually used.
\subsection{Volume fraction and interface}
Define $\alpha$ and the phase convention. Give the transport equation used in the code and discuss interface regularization.
\subsection{Initial and boundary conditions}
Document the implemented conditions: walls, open or closed top, initial state, and pressure gauge.

\section{Baseline PINN}
\subsection{Inputs, outputs, and normalization}
Describe nondimensional coordinates $(x^*,y^*,t^*)$, predicted variables $(u^*,v^*,p^*,\alpha)$, and reference scales.
\subsection{Network architecture}
Describe the baseline architecture in \texttt{src/model.py}: dense layers, Tanh activation, and sigmoid-bounded $\alpha$ output.
\subsection{Loss function}
Écrire la perte totale sous la forme
\[
\mathcal{L}=\lambda_c\mathcal{L}_{c}
+\lambda_\alpha\mathcal{L}_{\alpha}
+\lambda_m\mathcal{L}_{m}
+\lambda_0\mathcal{L}_{0}
+\lambda_w\mathcal{L}_{w}
+\lambda_p\mathcal{L}_{p}.
\]
Define each term according to the actual implementation, including units and normalization.
\subsection{Limitations of the baseline}
Describe observed interface diffusion or displacement, impact-region errors, and optimization difficulties. Distinguish measured observations from hypotheses.

\section{Proposed Methods}
\subsection{Optional Fourier feature encoding}
Présenter l'encodage des coordonnées, ses fréquences et le contrôle expérimental permettant de l'activer ou le désactiver.
\subsection{Optional adaptive sampling}
Décrire le bassin de candidats, le score combinant proximité de $\alpha=0.5$, gradient de $\alpha$ et résidus de continuité, transport et quantité de mouvement, ainsi que la fraction de points conservée uniformément.
\subsection{Temporal windows and weight transfer}
Describe window lengths, overlap, initialization from the previous window's weights, and replay strategy. Explain how temporal forgetting is mitigated.
\subsection{Optional adaptive loss weighting}
Décrire la méthode choisie (par exemple équilibrage fondé sur les normes de gradients), l'EMA éventuelle, les bornes des poids et la fréquence de mise à jour.
\subsection{Optional Adam followed by L-BFGS}
Présenter les étapes, le critère de passage à L-BFGS et l'utilisation de points fixes pendant chaque fermeture L-BFGS.
\subsection{Combinations and ablation study}
Préciser quelles options sont indépendamment activables et quelles combinaisons seront testées.

\section{Experimental Protocol}
\subsection{Test case and CFD reference}
Décrire le cas de ballottement, les propriétés des fluides, l'excitation, les conditions aux limites et les paramètres interFoam. Les données CFD ne doivent pas servir à l'entraînement si l'objectif est un PINN autonome ; préciser si elles sont utilisées uniquement pour l'évaluation.
\subsection{Experimental design}
Présenter au minimum : PINN de base, chaque option seule, combinaisons sélectionnées et répétitions avec graines différentes.
\subsection{Metrics}
\begin{itemize}[nosep]
\item erreur de fraction volumique et erreur de localisation de l'interface ;
\item distance entre interfaces (par exemple distance de Hausdorff ou distance moyenne), si applicable ;
\item erreur sur les vitesses et pression lorsque des références sont disponibles ;
\item résidus physiques et stabilité de l'entraînement ;
\item temps total, temps par époque, mémoire GPU/CPU et nombre d'évaluations ;
\item robustesse lors des impacts et de la rupture de vagues.
\end{itemize}
\subsection{Reproducibility details}
Report software versions, hardware, numerical precision, random seeds, compute budgets, hyperparameters, and scripts.

\section{Results}
\subsection{Implementation verification}
Présenter tests unitaires des résidus, normalisation, dérivées automatiques et options de configuration.
\subsection{Results de référence}
Présenter les performances du PINN non modifié.
\subsection{Effect of each improvement}
Une sous-section par option : Fourier, échantillonnage adaptatif, fenêtres temporelles, pondération adaptative et L-BFGS.
\subsection{Comparison of combined methods}
Présenter les interactions, les compromis précision/coût et les échecs éventuels.

\begin{table}[h!]
\centering
\caption{Comparaison des 16 configurations de modèle testées}
\begin{tabular}{lccccc}
\toprule
Configuration & RMSE u moy. & RMSE v moy. & RMSE p moy. & RMSE α moy. & Temps d'entraînement (s) \\
\midrule
model\_ff0\_as0\_ts0\_alw0 & 0.770 & 0.189 & 3724.733 & 0.490 & 45.083 \\
model\_ff0\_as0\_ts0\_alw1 & 0.774 & 0.188 & 3680.937 & 0.491 & 22.982 \\
model\_ff0\_as0\_ts1\_alw0 & 0.771 & 0.210 & 3641.053 & 0.492 & 13.475 \\
model\_ff0\_as0\_ts1\_alw1 & 0.768 & 0.199 & 3858.882 & 0.492 & 22.744 \\
model\_ff0\_as1\_ts0\_alw0 & 0.770 & 0.191 & 3876.398 & 0.491 & 14.481 \\
model\_ff0\_as1\_ts0\_alw1 & 0.772 & 0.188 & 3718.598 & 0.492 & 21.664 \\
model\_ff0\_as1\_ts1\_alw0 & 0.777 & 0.211 & 3537.408 & 0.493 & 16.883 \\
model\_ff0\_as1\_ts1\_alw1 & 0.777 & 0.195 & 3732.812 & 0.493 & 21.696 \\
model\_ff1\_as0\_ts0\_alw0 & 0.768 & 0.196 & 4207.534 & 0.486 & 17.317 \\
model\_ff1\_as0\_ts0\_alw1 & 0.766 & 0.197 & 4224.921 & 0.486 & 16.375 \\
model\_ff1\_as0\_ts1\_alw0 & 0.779 & 0.198 & 4169.999 & 0.481 & 6.685 \\
model\_ff1\_as0\_ts1\_alw1 & 0.777 & 0.210 & 4177.826 & 0.481 & 10.498 \\
model\_ff1\_as1\_ts0\_alw0 & 0.766 & 0.221 & 4198.084 & 0.484 & 8.281 \\
model\_ff1\_as1\_ts0\_alw1 & 0.767 & 0.231 & 4209.499 & 0.484 & 13.663 \\
model\_ff1\_as1\_ts1\_alw0 & 0.778 & 0.223 & 4139.737 & 0.480 & 8.122 \\
model\_ff1\_as1\_ts1\_alw1 & 0.778 & 0.265 & 4152.172 & 0.480 & 11.787 \\
\bottomrule
\end{tabular}
\label{tab:benchmark_results}
\end{table}

\textbf{Tableau 1 :} Résultats du benchmark sur 16 configurations (50 époques chacune). Les RMSE sont moyennes sur l'ensemble des temps d'évaluation. Le temps d'entraînement est donné en secondes.
\subsection{Interface and impact visualizations}
Insérer des champs $\alpha$, contours $\alpha=0.5$, champs de vitesse et profils temporels. Comparer à interFoam uniquement pour l'évaluation, si cette référence est disponible.

\section{Discussion}
\subsection{Pourquoi les modifications fonctionnent ou échouent}
Interpréter les résultats sans attribuer une causalité non démontrée.
\subsection{Computational cost and scalability}
Discuter le coût des dérivées d'ordre élevé, du score adaptatif et de L-BFGS.
\subsection{Limitations}
Décrire les limites du modèle physique, de la représentation de l'interface, des conditions aux limites et de la validation.
\subsection{Reproducibility and autonomy}
Préciser ce qui est appris sans données CFD et ce qui nécessite une référence externe pour l'évaluation.

\section{Conclusions and Future Work}
Summarize measured findings, the most effective improvements, and future work.

\section*{Code and Data Availability}
[URL du dépôt, commit utilisé, scripts, configurations et modalités de disponibilité des données.]

\section*{Author Contributions}
[À compléter.]

\section*{Funding and Conflicts of Interest}
[À compléter.]

\bibliographystyle{plain}
\bibliography{references}

\appendix
\section{Complete Parameters}
[Tableau des paramètres physiques, numériques et d'entraînement.]
\section{Detailed Loss Definitions}
[Formules exactes, conventions de signe et facteurs d'échelle.]
\section{Algorithms}
[Algorithms/pseudocode pour l'échantillonnage adaptatif et les fenêtres temporelles.]
\end{document}
